\section{Conclusion}
\label{sec:conclusion}

We set out to ask whether today's neoantigen immunogenicity tools can quantify the
\emph{magnitude} of a T-cell response, and to do so under a controlled,
apples-to-apples protocol rather than across incomparable papers. Three
conclusions follow. First, under a unified seven-step harmonization protocol,
existing tools do not quantify magnitude: across eight immunogenicity predictors
and one presentation proxy, every absolute Spearman correlation with the
continuous ELISpot readout falls below $0.33$, binary-discrimination differences
have broadly overlapping bootstrap intervals that shift with aggregation and
threshold, and no single tool is statistically distinguishable as best. Second, a
per-patient evaluation protocol recovers within-patient ranking ability that the
default global correlation conceals, exposing substantial individual variation
and arguing that benchmarks in this area should compute their metrics per patient
before aggregating. Third, response magnitude is a systematically overlooked gap,
but one bounded by a genuine biological ceiling: because the leading driver of
magnitude is a donor-specific precursor frequency invisible to peptide-and-allele
inputs, the explainable variance is structurally capped, and the value of
continuous prediction is best framed as ranking gain for clinical
top-$K$ selection rather than as high-precision regression.

Taken together, these findings reframe magnitude prediction as an open and
worthwhile---yet bounded---problem, and supply a reproducible benchmark and
per-patient protocol against which future continuous-magnitude predictors can be
measured. We hope the benchmark serves both as an honest baseline for what current
tools achieve and as a measuring stick for the quantitative immunogenicity models
that this gap invites.
